\subsection{Домашние задачи на десятую неделю.}
\begin{problem}
    Книга объёмом 500 страниц содержит 50 опечаток.
    Оцените вероятность того, что на случайно выбранной странице имеется не менее трёх опечаток.
    Используйте нормальное и пуассоновское приближения.
    Сравните результаты.
\end{problem}
\begin{solution} \leavevmode
    \begin{itemize}
        \item Беря в качестве $\lambda = 50/500 = 0.1$ мы получим, что $\P(\geq 3) = 1 - \sum\limits_{k = 0}^{2} \P( = 0) = 1 - \dfrac{e^{-0.1}}{0!} - \dfrac{0.1e^{-0.1}}{1!} - \dfrac{(0.1)^2 e^{0.1}}{2!} \approx 0.0002$.
    \end{itemize}
    Ну не будет же этой задачи на кр...
\end{solution}
\begin{problem}
    Случайные величины $\xi_n$ независимы в совокупности и принимают значения
    \begin{itemize}
        \item $\P(\xi_n = -2^n) = \P(\xi_n = 2^n) = 0.5$.
        \item $\P(\xi_n = -\sqrt {n}) = \P(\xi_n = \sqrt {n}) = 0.5$.
        \item $\P(\xi_n = -n) = 1/4, \P(\xi_n = n) = 1/4, \P(\xi_n = 0) = 1/2$.
        \item $\P(\xi_n = -n) = \dfrac{n + 1}{2n + 1}, \P(\xi_n = n + 1) = \dfrac{n}{2n + 1}$.
    \end{itemize}
    Выполнен ли для этих случайных величин ЗБЧ?
\end{problem}
\begin{solution}
    Для последовательности случайных величин $\{X_n\}$ с существующими первыми моментами известен критерий выполнимости ЗБЧ (Колмогорова):
    \[
        \Ex\biggr(\dfrac{\sum\limits_{k = 1}^{n}(X_k - \Ex X_k)^2}{n^2 + \sum\limits_{k = 1}^n(X_k - \Ex X_k)^2}\biggr) \ra 0, n \ra +\infty.
    \]
    \begin{itemize}
        \item Посчитаем математические ожидания и дисперсии величин из последовательности: $\Ex \xi_n = 0, \D \xi_n = 4^n$ (и, при этом, $\xi_n^2$ является $4^n$ с вероятностью 1).
        Тогда
        \[
            \dfrac{\sum\limits_{k = 1}^n 4^k}{n^2 + \sum\limits_{k = 1}^{n}4^k} \ra 1, n \ra +\infty.
        \]
        А значит для этой последовательности не выполнен ЗБЧ.
        \item Посчитаем математические ожидания и дисперсии величин из последовательности: $\Ex \xi_n = 0, \D \xi_n = n$ (и, при этом, $\xi_n^2$ -- это $n$ с вероятностью 1).
        Тогда
        \[
            \dfrac{\sum\limits_{k = 1}^n n}{n^2 + \sum\limits_{k = 1}^n n} = \dfrac{\dfrac{n(n + 1)}{2}}{n^2 + \dfrac{n(n + 1)}{2}} = \dfrac{n^2 + n}{3n^2 + n} \ra 1/3, n \ra +\infty.
        \]
        А значит для этой последовательности также не выполнен ЗБЧ.
        \item Посчитаем математические ожидания и дисперсии для этой последовательности: $\Ex \xi_n = 0, \D \xi_n = \Ex \xi_n^2 = n^2/2$.
        Поскольку $\sum\limits_{k = 1}^n (X_k - \Ex X_k)^2 \leq \sum\limits_{k = 1}^{n}k^2 = \dfrac{n(n + 1)(2 + 1)}{6}$ с вероятностью $1$, то
        \[
            \Ex \biggr(\dfrac{\sum\limits_{k = 1}^n (X_k - \Ex X_k)^2}{n^2 + \sum\limits_{k = 1}^{n} (X_k - \Ex X_k)^2} \biggr) \geq \dfrac{\Ex \sum\limits_{k = 1}^n (X_k - \Ex X_k)^2}{n^2 + \dfrac{n(n + 1)(2n + 1)}{6}} \ra \dfrac{1}{2}, n \ra +\infty.
        \]
        А значит для этой последовательности не выполнен ЗБЧ.
        \item Посчитаем математические ожидания и дисперсии величин в последовательности:
        \[\Ex \xi_n = 0, \D \xi_n = \dfrac{n^2(n + 1) + (n + 1)^2 n}{(2n + 1)} = \dfrac{n(n + 1)(2n + 1)}{(2n + 1)} = n(n + 1).\]
         Аналогично предыдущему пункту оценим подынтегральную функцию и в пределе придём не к нулю, а следовательно ЗБЧ для этой последовательности будет невыполнен.
    \end{itemize}
\end{solution}
\begin{problem}
    При каких $\alpha > 0$ к последовательности случайных величин $\{\xi_n\}_{n \in \N}$ применим ЗБЧ, если $\P(\xi_n = -n\alpha) = \P(\xi_n = n \alpha) = 2^{-n}$, а $\P(\xi_n = 0) = 1 - 2^{1 - n}$.
\end{problem}
\begin{solution}
    При любых $n$ и $\alpha$ справедливо $\Ex \xi_n = 0$.
    Дисперсии выражаются как
    \[
        \D \xi_n = (n\alpha)^2 \dfrac{2}{2^n} = \alpha^2 n^{2}2^{1 - n}.
    \]
    Поскольку
    \[
        0 \leq \Ex \biggr(\dfrac{\sum\limits_{k = 1}^n \xi_k^2}{n^2 + \sum\limits_{k = 1}^n \xi_k^2}\biggr) \leq \dfrac{\Ex\sum\limits_{k = 1}^n \xi_k^2}{n^2} = \dfrac{\sum\limits_{k = 1}^n (k\alpha)^2 2^{1 - n}}{n^2} = \dfrac{2\alpha^2 n(n + 1)(2n + 1)}{6n^2 2^{n}} \ra 0, n \ra +\infty.
    \]
    А значит для этой последовательности выполнен ЗБЧ при любом $\alpha$.
\end{solution}
\begin{problem}
    Учебник издан тиражом 100000 экземпляров.
    Вероятность того, что учебник сброшюрован неправильно, равна 0,0001.
    Оцените с помощью подходящей предельной теоремы вероятность того, что тираж содержит не менее шести бракованных книг.
    Сравните с точным значением.
\end{problem}
\begin{solution}
    Пусть $S_n = \sum\limits_{k = 1}^{100000} Be(0.0001) \sim Binom(100000, 0.0001)$. \\
    Используем Пуассоновское приближение с $\lambda = 100000 \cdot 0.0001 = 10$:
    \[
        \P(S_n \geq 6) \sim 1 - \sum\limits_{k = 1}^{5}\dfrac{\lambda^k e^{-\lambda}}{k!} \sim 0.9329140
    \]
    При точном значении в $\P(S_n \geq 6) = 1 - \sum\limits_{k = 1}^{5}\binom{10^5}{k}p^{k}q^{10^5-k} \approx 0.9329235$.
\end{solution}
\begin{problem}
    Вероятность рождения мальчика равна 0,51.
    Оцените с помощью подходящей предельной теоремы вероятность того, что среди 1000 новорожденных окажется ровно 500 мальчиков.
    Сравните с точным значением.
\end{problem}
\begin{solution}
    Пусть $S_n = \sum\limits_{k = 1}^{1000} Be(0.51) \sim Binom(1000, 0.51)$. \\
    Тогда по локальной теореме Муавра-Лапласа:
    \[
        \P(S_n = 500) \sim \dfrac{1}{\sqrt{2\pi n p q}}e^{-x^2/2}, \quad x = \dfrac{k - np}{\sqrt{npq}}.
    \]
    Численно это примерно выражается как $\P(S_n = 500) \sim 0.02066$. \\
    Точный подсчёт даёт такую оценку:
    \[
        \P(S_n = 500) = \binom{1000}{500} 0.51^{500}0.49^{500} \approx 0.02065.
    \]
\end{solution}
\begin{problem}
    Пусть множество элементарных исходов некоторого вероятностного пространства счётно.
    Следует ли в этом случае из сходимости случайных величин по вероятности сходимость почти наверное?
\end{problem}
\begin{solution}
    Заметим, что мы можем считать, что $\forall w \in \Omega \hookrightarrow \P(w) > 0$, поскольку точки меры нуль мы можем просто выкинуть и это никак не повлияет на сходимость (и у нас счётное вероятностное пространство).
    Пусть $w \in \Omega: \P(w) = \alpha \in (0, 1)$.
    По определению сходимости по вероятности $\P(|X_n - X| > \alpha/2) \ra 0, n \ra +\infty$.
    Но, если $X_n(w) \not\ra X(w)$, то $\P(|X_n - X| > \alpha) \geq \alpha \not\ra 0, n \ra +\infty$ -- следовательно у нас есть поточечная сходимость во всякой точке положительной меры.
    А значит в этом случае из сходимости по мере следует сходимость почти наверное.
\end{solution}